\documentclass[11pt]{letter}
\usepackage[utf8]{inputenc}
\usepackage[T1]{fontenc}
\usepackage[margin=1in]{geometry}
\usepackage{hyperref}
\hypersetup{colorlinks=true, urlcolor=blue}

\signature{Hemanth Manchabale Papachappa \\ Aniriuddha Ganguly}
\address{}
\date{September 2026}

\begin{document}

\begin{letter}{
  Editor-in-Chief \\
  IEEE Journal of Translational Engineering in Health and Medicine (JTEHM) \\
  IEEE Engineering in Medicine and Biology Society
}

\opening{Dear Editor-in-Chief,}

We submit for your consideration the manuscript \textbf{``Geodesic Diagnostic
Trajectories (GDT): A Differential Geometry Framework for Adaptive Decision
Support and Operational Optimization in EHR Systems''} by Hemanth Manchabale
Papachappa and Aniriuddha Ganguly, for review as a regular article in
\textit{IEEE Journal of Translational Engineering in Health and Medicine}.

\textbf{Clinical problem and translational motivation.} Chart review in electronic
health record (EHR) systems is a multi-pivot, cognitively demanding task. A
physician tracking a critically ill patient may cycle through culture results,
renal dosing adjustments, and anticoagulation orders within a single session---each
pivot rendering the prior retrieval context a source of interference rather than
assistance. Standard session-based retrievers accumulate this stale context without
detection, degrading result relevance at the moment clinical decision speed matters
most. GDT is our direct engineering response to this failure mode, designed to move
from academic benchmark toward real-time clinical deployment.

\textbf{Technical contribution.} GDT encodes the clinician's evolving diagnostic
intent as a trajectory on the Riemannian manifold $\mathcal{S}_d^+$ of symmetric
positive definite (SPD) matrices. A geodesic shift ratio $\kappa_t$ quantifies
abrupt topic pivots relative to recent trajectory pace; a soft gate driven by
$\kappa_t$ suppresses stale context with suppression strength bounded by a formally
proved proposition (Proposition~1). A JEPA-style latent predictor anticipates the
next query and populates a prefetch cache; prefetch volume is scaled by trajectory
confidence so that uncertain post-pivot forecasts do not contaminate the cache.
An M/M/c queueing model connects the resulting latency gains to ward-level
operational throughput estimates.

\textbf{Clinical benchmark and key findings.} We evaluated GDT against CMA, BM25,
and TF-IDF in a randomized four-arm within-subject crossover study spanning 124
sessions on 31 multi-pivot clinical vignettes from MIMIC-III/IV, eICU, and HiRID.
GDT reduced mean retrieval latency by 26.9\% relative to TF-IDF control
($p{<}0.001$, $d{=}{-4.12}$). Proposition~1 held across all 1,074 gate-triggered
transitions without exception. The M/M/c projection estimates a 4.1-minute
reduction in per-reviewer patient processing time during peak-acuity shifts.
Equally important is the finding that all four conditions achieved 0\% task
completion---an honest diagnostic result that isolates the vocabulary mismatch
inherent in sparse lexical retrieval and identifies dense retrieval integration as
the necessary next step.

\textbf{Fit for JTEHM.} This manuscript aligns directly with JTEHM's mission of
translational engineering for health and medicine. The contributions span three
translational dimensions: (1) a theoretical framework with a proved interference
bound---the kind of formal guarantee that regulatory bodies and clinical governance
committees increasingly expect from deployed CDS; (2) a simulation benchmark that
makes the system's limitations explicit rather than obscuring them; and (3) an
operational analytics model that translates sub-second engineering improvements
into workflow metrics meaningful to hospital administrators. The paper also
discusses concrete deployment prerequisites---SMART on FHIR integration,
NASA-TLX human-subject validation, fairness auditing---situating it firmly in the
translational rather than purely algorithmic literature.

\textbf{Confirmations.}
\begin{itemize}
  \item This work has not been published and is not under consideration elsewhere.
  \item All authors have approved the manuscript and agree with the submission.
  \item All datasets used (MIMIC-III/IV, eICU, HiRID) are publicly available and
        de-identified; no new protected health information was collected.
  \item The evaluation pipeline will be shared on GitHub upon acceptance.
  \item There are no conflicts of interest to disclose (see accompanying statement).
\end{itemize}

We believe GDT's combination of formal theoretical foundations, transparent
reporting of null retrieval-quality findings, and operational throughput modeling
makes it a timely and relevant contribution to JTEHM's readership. We look forward
to the editors' and reviewers' feedback.

\closing{Sincerely,}

\end{letter}
\end{document}
